\section{Related Works and Motivation}
This section reviews poisoning attacks on \emph{RAG systems} and GSEs, summarizes key evaluation criteria, and highlights challenges unique to GSEs.

\label{sec:related}
\subsection{Vulnerabilities of GSEs}
We introduce a poisoning attack in GSEs.
Zou et al.~\cite{poisonedrag-arxiv24} reveal and formalize an attack method called \emph{PoisonedRAG} that exploits \emph{RAG systems} to generate attacker-intended answers by placing malicious content into external databases.
Even if attackers can inject a small volume of malicious text into external databases, \emph{RAG systems} can produce attacker-intended answers for specific questions.
This malicious content includes text that attackers want displayed as answers to specific questions, especially false information such as ``\textit{OpenAI's CEO is Tim Cook.}''
PoisonedRAG attacks target closed-ended questions like ``\textit{Who is the CEO of OpenAI?}'' rather than open-ended questions like ``\textit{What are the latest trends in AI?}''~\cite{poisonedrag-arxiv24}.

In fact, several studies support the effectiveness of this attack.
Chen et al.~\cite{chen2022rich} show that in \emph{RAG systems}, when content containing two different claims (e.g., \textit{OpenAI's CEO is Tim Cook.} and \textit{OpenAI's CEO is Sam Altman.}) is placed and answers are generated for each claim, the answers may include each respective claim.
Wang et al.~\cite{wang2025ragconflictevidence} and Liu et al.~\cite{liu2025conflicts} demonstrate that even when there is an imbalance in the volume of supporting documents for each claim, \emph{RAG systems} tend to favor the answer with more supporting evidence, yet the underrepresented claim can still appear in responses.
These studies support the effectiveness of PoisonedRAG attacks and demonstrate that RAG may cite malicious content even when attackers place only a small volume of it.

To succeed in a poisoning attack, the malicious content must satisfy two conditions: \emph{Retrieval Condition} and \emph{Generation Condition}.
Retrieval Condition means that malicious content (such as text or HTML placed by attackers) is selected as top-\( k \) relevant content through retrieval tasks for target questions.
In GSEs, Retrieval Condition means that content on the web is retrieved as web pages ranked highly in search engine results by SE.
Generation Condition means that the content is used as reasoning knowledge in the LLM and reflected in the generated answers.
Aggarwal et al.~\cite{geo-apmr+24} propose \emph{generative engine optimization} (GEO), a method that optimizes website text structure and content from the perspective of Generation Condition to increase the exposure of web page content as citations in answers when the website passes the Retrieval Conditions.
GEO theoretically supports that GSEs are vulnerable to poisoning attacks.

PoisonedRAGは政治分野における有権者の認知操作と直接結びつく攻撃手法である．
民主主義国家を対象とした多数の事例研究は，Web上に意図的に配置された
虚偽・誤導情報が選挙や政治的議論に影響を及ぼしてきたことを示している．

米国大統領選挙では，ロシアによる大規模な情報操作工作が
複数の事例研究で確認されてきた~\cite{badawy2018analyzing,zannettou2018disinformation,cinus2025cross}．
Badawy et al.~\cite{badawy2018analyzing}とZannettou et al.~\cite{zannettou2018disinformation}は，
2016年大統領選挙期間中のSNSにおける発信を分析し，
ロシア系trollアカウントによる偽装活動が
複数のSNSプラットフォームで再発信されていたことを示した．
同様の偽装活動は2024年米国大統領選挙においても継続的に観測されており，
Cinus et al.~\cite{cinus2025cross}は，
複数のSNSプラットフォームを横断する協調的な偽装活動を検出し，
ロシア系メディアがプラットフォームを跨いで組織的に拡散されていたことを明らかにした．

同様の脅威は米国以外の民主主義国家においても観測されている．
日本においては，2024年衆院選・2025年参院選において，
日本国外からのtrollやbotアカウントによってSNSプラットフォームに発信された情報が
特定政党の主張を増幅し，世論形成に影響を与えている~\cite{nikkei2025sanseito,nippon2025russian}．
さらに2025年には，中国当局と関係するアカウントがLLMを利用して
政治家への組織的な誹謗キャンペーンを計画していたことが
OpenAIの公式レポートで明らかにされている~\cite{openai2026disrupting}．

以上のように，攻撃者がWeb上にコンテンツを意図的に配置することで
GSEを通じて有権者の認知に影響を与える行為は，
特定の国を問わず世界的に継続的に発生している現実の脅威である．

\subsection{Citation Evaluation}
\label{subsec:citation_evaluation}
We first introduce seminal studies that aim to evaluate \emph{RAG systems}.
These studies highlight that in ideal \emph{RAG systems}, all claimed sentences in an answer output from the \emph{RAG system} are supported by external sources from databases as \emph{citations} (citation coverage) and the sentences with citations accurately and faithfully reflect the claims (reflection accuracy).
To analyze them, numerous studies evaluate citation coverage and reflection accuracy, specifically focusing on \emph{RAG systems}~\cite{citeeval-acl25, ragas-eacl24, ares-naacl24, evaluating-verifiability-in-generative-search-engines-emnlp23, enabling-llms-to-generate-text-with-citations-emnlp23, ranking-generated-summaries-by-correctness-acl2019, evaluating-the-factual-consistency-of-abstractive-text-summarization-emnlp20, alignscore-acl23, feqa-acl20}.

A core component of the analysis of reflection accuracy is faithfulness evaluation metrics~\cite{finegrainedcitationevaluationgenerated-arxiv24}.
Faithfulness evaluation metrics are typically classified into three methods: entailment-based~\cite{ranking-generated-summaries-by-correctness-acl2019, true-reevaluating-factual-consistency-arxiv22, evaluating-the-factual-consistency-of-abstractive-text-summarization-emnlp20, alignscore-acl23}, similarity-based~\cite{bertscore-arxiv20, bartscore-neurips21}, and QA-based~\cite{feqa-acl20}.
Similarity-based methods are most appropriate for calculating reflection accuracy.
Similarity-based methods quantitatively measure semantic similarity between two texts and rely on neural encoder models.
Zhang et al.\ show that similarity-based methods achieve higher precision than entailment-based and QA-based methods in tasks evaluating consistency between long documents and their summaries~\cite{finegrainedcitationevaluationgenerated-arxiv24}.
This approach aligns with GES behavior of summarizing web page content retrieved from search results and generating answers with citations.
Citation evaluation applies these faithfulness evaluation methods to each sentence \( l_i \) of generated answer \( r \) with citations \( C_i \), and evaluates reflection accuracy for the entire answer \( r \).

However, they do not examine which publishers of the content are selected as citations and assess the authority and trustworthiness of those sources.
Thus, a gap remains between the seminal studies on \emph{RAG systems} and their applicability to GSEs.
In \emph{RAG systems}, poisoning attack studies focus on how much malicious content attackers need to inject into external databases to succeed in the poisoning attack, while since GSEs retrieve information from the open web where attackers can inject malicious content more easily than in \emph{RAG systems}, 
Therefore, the main concern for GSEs is whether they retrieve and incorporate even a small amount of malicious content injected by attackers during web searches, since this can enable a successful poisoning attack as previously described. 
Thus, existing evaluation methods alone are insufficient to assess how vulnerable GSEs are to poisoning attacks or to fully characterize their attack surface.

\subsection{Authority and Trustworthiness}
Several studies propose methods to improve the inclusion of authoritative sources in citations generated by \emph{RAG systems} and GSEs.
TrustRAG~\cite{zhou2025trustragenhancingrobustnesstrustworthiness} identifies and removes malicious documents during retrieval by clustering candidate documents and matching them with the internal knowledge of the LLM.
AuthGR~\cite{lee2026relevanceauthorityauthorityawaregenerative} computes an authority score for web pages using a vision-language model that evaluates both the textual content of the page and visual features extracted from page screenshots, such as layout quality, design, and advertisement density. The system then preferentially cites web pages with high authority scores.

\subsection{Authority and Trustworthiness}
Several studies propose methods to improve the inclusion of authoritative sources in citations generated by \emph{RAG systems} and GSEs.
TrustRAG~\cite{zhou2025trustragenhancingrobustnesstrustworthiness} identifies and removes malicious documents during retrieval by clustering candidate documents and matching them with the internal knowledge of the LLM.
AuthGR~\cite{lee2026relevanceauthorityauthorityawaregenerative} computes an authority score for web pages using a vision-language model that evaluates both the textual content of the page and visual features extracted from page screenshots, such as layout quality, design, and advertisement density. The system then preferentially cites web pages with high authority scores.

However, evaluating only the content of citations is insufficient for assessing citation safety under poisoning attacks.
As described in Section~\ref{subsec:citation_evaluation}, because anyone can publish content on the web, attackers can deliberately optimize the scores of content-based evaluation metrics while still injecting malicious information into the web.
実際，EU DisinfoLabが報告した\emph{Doppelganger}作戦~\cite{eudisinfolab2022doppelganger,foreignaffairs2024doppelganger}では，
正規メディアのドメイン名・デザインを模倣した
700以上の偽造ニュースサイトが公開され，
正規メディアに偽装した形で配信されていたことが示されている．
攻撃者がWeb上に「正規メディアに見える」コンテンツを大量配置するという手法は，
コンテンツ自体を評価するアプローチを無効化する．
Therefore, for GSEs, we need to evaluate not only the content of citations, but also the authority and trustworthiness of their publishers.

Furthermore, they typically assume the presence of poisoning attacks on \emph{RAG systems} and primarily focus on developing defenses or system improvements. 
However, they do not provide a systematic understanding or measurement of the actual extent of citation selection bias and vulnerability found in production GSEs.


\subsection{Personalization of LLMs}
Personalization is a key factor in improving the user experience of chat-based applications such as ChatGPT, Claude, and Gemini.
In personalization, the general approach is to embed the user's attribute information together with the user's question.
This approach enables the applications to generate responses that more accurately reflect the user's preferences, background, and needs.

Several studies investigate how personalization affects answers generated from the internal knowledge of LLMs.
Tornberg et al.~\cite{political-bias-audit-arxiv-2026} analyze political ideology bias in representative and widely used LLMs.
The study shows that LLMs produce left-leaning responses by default and LLMs shift strongly to the right for conservative users but only slightly to the left for progressive users, showing an asymmetric response to personalization.

This effect is not limited to political topics.
Kearney et al.~\cite{kearney2025languagemodelschangefacts} show that across several domains, answers change systematically depending on user attributes such as race, gender, and age.
The influence of user attributes also extends to tasks that have objectively correct answers.
Vijjini et al.~\cite{exploreing-safery-in-personalization-arxiv2026} show that LLM response performance varies systematically even in tasks such as mathematical reasoning, general knowledge, and programming.

However, these previous studies focus on biases that arise from the internal knowledge of LLMs and do not examine GSEs, where answers are generated using web search as a broad external context. 
To the best of our knowledge, no prior work has examined how personalization affects the citation behavior of GSEs, including aspects such as source authority and reflection accuracy.

\subsection{Research Motivation}
To better understand the attack surface of GESs, our goal is to provide a novel evaluation framework. 
Our framework introduces new evaluation criteria centered on the publisher attributes of citations in generated answers, enabling a systematic assessment of how robust GESs are to poisoning attacks.

Specifically, our framework addresses the following questions: 
(1) Which publishers' content do GESs preferentially cite, and does this reveal any bias in citation selection from the perspective of publishers? 
(2) Which publishers’ content has the power to influence the faithfulness and accuracy of the generated answers?
(3) How does personalization impact citation behavior in GESs?

Importantly, our primary focus is on establishing robust evaluation criteria for assessing the vulnerability of GESs to poisoning attacks.
We do not aim to demonstrate practical attacks or identify actual instances of malicious content injection.
Rather, our goal is to assess the potential risk of poisoning attacks to reveal the underlying attack surfaces, and to clarify the degree of resilience that GESs exhibit against such attacks.

本研究では、我々のフレームワークを政治領域に適用し、GSEが情報伝達の基盤となる社会において、民主主義の根幹である政治情報の伝達に関するリスクを、GSEの引用行動を通じて明らかにすることを目指す。
特に、有権者が特定の政党について情報収集を行うタスクに着目し、政党別・イデオロギー別・国別、ならびに有権者のペルソナ（政治的知識やイデオロギー）の違いによって、frameworkを通じたattack surfaceを特定する。
そのattack surfaceを踏まえ、どの領域において、攻撃者が政治情報を意図的に湾曲させ、情報伝達の偏りや歪みが生じる可能性があるかを考察する。

